\section{A Taxonomy of Jailbreak Judge Failure Modes}

We analyze the conflict cases collected by TVD-Conflict and organize them into
six recurring failure modes. These categories describe the source of the
judgment error rather than the surface form of the jailbreak.
Table~\ref{tab:failure-mode-prevalence} will report their prevalence after the
conflict set has been annotated.

\begin{table*}[t]
  \centering
  \caption{Occurrence of recurring judge vulnerabilities in conflicts obtained
  from attack--victim outputs ($n=717$) and \textsc{TVD}-Conflict ($n=1{,}681$,
  after excluding 21 judge-output artifacts). The first four modes are
  \emph{primary} attributions --- every conflict is assigned to exactly one,
  so each column sums to 100\%. The last two modes are \emph{cross-cutting}
  attributes that can co-occur with any primary mode: $^{\dagger}$share of all
  conflicts; $^{\ddagger}$share of binary-judge conflicts only (408
  attack-derived, 771 \textsc{TVD}-Conflict), since this mode is defined on
  the binary judges' structured sub-labels. All counts are exact: the
  hallucinated-evidence counts come from manually verifying every
  programmatically flagged absence claim against the full response text, and
  the internally-incoherent counts are computed mechanically from the
  sub-labels.}
  \label{tab:failure-mode-prevalence}
  \small
  \setlength{\tabcolsep}{7pt}
  \begin{tabular}{l rr rr}
    \toprule
    & \multicolumn{2}{c}{Attack-Derived Conflicts}
    & \multicolumn{2}{c}{\textsc{TVD}-Conflict} \\
    \cmidrule(lr){2-3}\cmidrule(lr){4-5}
    Vulnerability & Count & Share (\%) & Count & Share (\%) \\
    \midrule
    Criterion operationalization divergence & 349 & 48.7 & 649 & 38.6 \\
    Refusal--substance inconsistency & 65 & 9.1 & 405 & 24.1 \\
    Salience misallocation & 248 & 34.6 & 506 & 30.1 \\
    Representation/reasoning limitation & 55 & 7.7 & 121 & 7.2 \\
    \midrule
    Hallucinated evidence$^{\dagger}$ & 41 & 5.7 & 99 & 5.9 \\
    Internally incoherent verdict$^{\ddagger}$ & 112 & 27.5 & 107 & 13.9 \\
    \bottomrule
  \end{tabular}
\end{table*}


\paragraph{Criteria ambiguity.}
The judges operationalize the stated safety criterion differently, such as
treating harmfulness, policy violation, and attack success as interchangeable.

\paragraph{Refusal but harmful.}
A response begins with refusal language but subsequently provides information
that materially advances the harmful request. Some judges anchor on the refusal
marker, while others evaluate the supplied content.

\paragraph{Compliance but safe.}
A response adopts the form of compliance without delivering harmful or usable
information. Judges that rely on affirmative framing can mistake stylistic
compliance for harmful assistance.

\paragraph{The 50--50 problem.}
The available evidence supports two plausible readings of comparable strength,
causing verdicts to split across judges or rubric formulations.

\paragraph{Distraction.}
An evaluator focuses on wrappers, disclaimers, fictional settings, mathematical
notation, or peripheral detail instead of the behavior that determines safety.

\paragraph{Insufficient capability.}
The judge fails to decode, follow, or reason over the representation required to
identify the harmful behavior, including obfuscation and multi-step context.
